\documentclass[12pt]{article}

\usepackage{amssymb,amsmath,amsthm}

% gaya teorema (saya menyukai semua lingkungan memakai pencacah yang sama); adaptasi berlisensi CC BY 4.0
\theoremstyle{plain}
\newtheorem{theorem}{Teorema}[section]
\newtheorem{lemma}[theorem]{Lema}
\newtheorem{proposition}[theorem]{Proposisi}
\newtheorem{corollary}[theorem]{Korolari}
\newtheorem{conjecture}[theorem]{Konjektur}
\newtheorem{exercise}[theorem]{Latihan}
\theoremstyle{definition}
\newtheorem{definition}[theorem]{Definisi}
\newtheorem{example}[theorem]{Contoh}
\theoremstyle{remark}
\newtheorem{remark}[theorem]{Catatan}

% beberapa pengaturan penomoran
\numberwithin{equation}{section}

\usepackage{fancyhdr}
\usepackage{lastpage}
\setlength{\headheight}{15.2pt}
\renewcommand{\footrulewidth}{0.4pt}% nilai bakunya 0pt
\setlength{\footskip}{30pt}
\pagestyle{fancy}

\makeatletter
\let\ps@plain\ps@fancy 
\makeatother

\lhead{MATH 742}
\chead{Representasi atas medan yang tidak tertutup}
\rhead{Musim Semi 2023}
\lfoot{Revisi Terakhir: \today}
\cfoot{}
\rfoot{\thepage\ dari \pageref{LastPage} }

\begin{document}

%%%%%%%%%%%%%%%%%%
%%%%%%%%%%%%%%%%%%
\title{Representasi atas medan yang tidak tertutup}
\author{Alexander Duncan}

\maketitle

Di seluruh catatan ini, $k$ adalah medan berkarakteristik $0$
dan $\overline{k}$ adalah penutupan aljabar dari $k$.
(Bagi pembaca yang kurang nyaman dengan perluasan Galois tak hingga,
$\overline{k}$ dapat diganti dengan penutupan Galois
dari sejumlah hingga perluasan medan hingga yang relevan
dan muncul dalam konteks-konteks yang dibahas di bawah.)

Sebagian besar pembahasan ini terinspirasi oleh \cite[\S{12--13}]{Serre}.

\section{Gelanggang Representasi melalui Modul}

Misalkan $G$ suatu grup hingga.

Ingat bahwa $kG$ adalah aljabar grup $G$ atas $k$.
Karena $k$ berkarakteristik $0$, kita tahu bahwa $kG$ semisederhana.
Kita memiliki dekomposisi kanonik
\[
kG \cong \bigoplus_{i=1}^r A_i
\]
dengan $A_1,\ldots,A_r$ aljabar-aljabar $k$ sederhana.

Setiap $A_i$ merupakan aljabar sederhana sentral atas $F_i$
untuk suatu perluasan medan hingga $F_i/k$.
Misalkan $n_i = \deg(A_i)$ adalah \emph{derajat} dari $A_i$.
Secara ekuivalen, berlaku $n_i^2 = \dim_{F_i}(A_i)$
atau
\[
A_i \otimes_{F_i} \overline{k} \cong
\operatorname{M}_{n_i}(\overline{k})
\]
sebagai aljabar-aljabar $\overline{k}$.

Lebih tepatnya, terdapat suatu aljabar-$F_i$ pembagian sentral $D_i$
dan suatu bilangan bulat positif $m_i$ sedemikian sehingga
\[
A_i \cong \operatorname{M}_{m_i}(D_i)
\]
sebagai aljabar-aljabar $F_i$.
Baik $m_i$ maupun kelas isomorfisme $D_i$ ditentukan secara tunggal.
Misalkan $d_i := \operatorname{ind}(A_i)$ menyatakan indeks Schur dari $A_i$.
Kita mempunyai $d_i = \deg(D_i)$, yang ekuivalen dengan
$d_i^2=\dim_{F_i}(D_i)$.  Selain itu, $n_i=m_id_i$.

Setiap $A_i$ mempunyai modul sederhana $V_i$, yang tunggal hingga isomorfisme.
Karena $V_i \cong D_i^{\oplus n_i}$, kita memperoleh
\[
\dim_k V_i = [F_i:k]m_id_i^2 = [F_i:k]d_in_i .
\]
Ruang-ruang vektor $V_1,\ldots,V_r$ tepat merupakan
representasi tak tereduksi dari $G$ atas $k$.
\emph{Indeks Schur} dari $V_i$ didefinisikan sebagai
indeks Schur $d_i$ dari aljabar terkait $A_i$.

Misalkan $\Gamma$ adalah grup Galois absolut
$\operatorname{Gal}(\overline{k}/k)$.
Misalkan $X_i :=
\operatorname{Hom}_{k-\operatorname{alg}}(F_i,\overline{k})$.
Perhatikan bahwa $X_i$ memiliki aksi kiri oleh $\Gamma$ melalui
pascakomposisi.
Secara informal, $X_i$ dapat dipandang sebagai himpunan akar polinomial
minimal dari suatu elemen primitif bagi $F_i/k$.
Kita mempunyai isomorfisme aljabar-$\overline{k}$
\[
\omega : F_i \otimes_k \overline{k} \cong
\bigoplus_{\sigma \in X_i} \overline{k}
\]
melalui $\omega(f \otimes \ell)_{\sigma} := \sigma(f)\ell$.

Untuk setiap $\sigma \in X_i$, kita dapat mendefinisikan hasil kali tensor
$A_i \otimes_\sigma \overline{k}$ dari aljabar-aljabar $F_i$
dengan memakai pemetaan struktur $F_i \to A_i$ di sebelah kiri dan
homomorfisme $\sigma : F_i \to \overline{k}$ di sebelah kanan.
Melalui rangkaian isomorfisme aljabar-$\overline{k}$
\[
A_i \otimes_k \overline{k} \cong 
A_i \otimes_{F_i} F_i \otimes_k \overline{k}
\cong A_i \otimes_{F_i} \bigoplus_{\sigma \in X_i} \overline{k}
\]
kita memperoleh isomorfisme
\[
\Omega : A_i \otimes_k \overline{k} \cong
\bigoplus_{\sigma \in X_i} A_i \otimes_\sigma \overline{k} .
\]

Misalkan $W_{i,\sigma}$ adalah representasi tak tereduksi
dari $\overline{k}G$ yang bersesuaian dengan $A_i \otimes_\sigma \overline{k}$.
Dengan demikian, kita memperoleh dekomposisi berikut
\[
V_i \otimes_k \overline{k} \cong
\bigoplus_{\sigma \in X_i} W_{i,\sigma}^{\oplus d_i}
\]
untuk setiap representasi tak tereduksi $V_i$ dari $kG$.

Jika $\rho_{i,\sigma} : G \to \operatorname{GL}(W_{i,\sigma})$ adalah
pemetaan-pemetaan yang bersesuaian, perhatikan bahwa
\[
\tau \circ \rho_{i,\sigma} = \rho_{i,\tau(\sigma)}
\]
untuk setiap $\tau \in \Gamma$.
Jadi terdapat aksi kiri $\Gamma$ pada himpunan representasi tak tereduksi
dari $G$ atas $\overline{k}$
dengan orbit-orbit yang bersesuaian dengan himpunan-$\Gamma$ $X_i$.

\section{Teori Karakter}

Sejauh ini kita baru mengembangkan teori karakter atas $\mathbb{C}$.
Namun, karena kita membahas representasi berdimensi hingga dari
grup hingga, berdasarkan prinsip Lefschetz
semua hasil yang telah dibahas berlaku atas sembarang medan tertutup secara aljabar
yang berkarakteristik $0$.
Khususnya, kita dapat mengidentifikasi gelanggang-gelanggang representasi
$R(G)=R_{\mathbb{C}}(G)=R_{\overline{k}}(G)$
dengan $\overline{k}$ penutupan aljabar dari $k$.

Ingat bahwa gelanggang representasi $R_k(G)$ dari $G$ atas $k$
adalah grup aditif representasi virtual atas $k$
dengan perkalian yang diperoleh dengan memperluas hasil kali tensor
pada subsemigelanggang $R_k^+(G)$.
Dari pembahasan di atas, $R_k(G)$ adalah grup abelian bebas
dengan basis $[V_1],\ldots,[V_r]$.

Pemetaan $[V] \mapsto [V \otimes_k \overline{k}]$
memberikan homomorfisme gelanggang injektif yang kanonik
\[
R_k(G) \hookrightarrow R_{\overline{k}} = R(G)
\]
dan karenanya homomorfisme gelanggang injektif kanonik
\[
R_k(G) \hookrightarrow C(G)
\]
dengan $C(G)$ himpunan fungsi kelas kompleks pada $G$.

Seperti dalam kasus kompleks, untuk suatu representasi $(V,\rho)$
dari $G$ atas $k$, kita mendefinisikan karakter
$\chi_V : G \to k$
melalui teras $\chi_V(g) := \operatorname{tr}(\rho(g))$.
Definisi ini sesuai dengan pemetaan $R_k(G) \hookrightarrow C(G)$ di atas.
Selanjutnya, kita mengidentifikasi $R_k(G)$ dengan subgelanggang dari $R(G)$ dan $C(G)$
dalam pembahasan berikut.

Misalkan $\chi_1,\ldots,\chi_r$ adalah karakter dari representasi tak tereduksi $G$
atas $k$ yang bersesuaian dengan aljabar $A_1,\ldots, A_r$ di atas.
Misalkan $\psi_{i,\sigma} : G \to \overline{k}$ adalah karakter
yang bersesuaian dengan representasi tak tereduksi $W_{i,\sigma}$
yang didefinisikan pada bagian sebelumnya.

\begin{theorem}
Untuk setiap $1 \le i \le r$, berlaku
\[
\chi_i = d_i \sum_{\sigma \in X_i} \psi_{i,\sigma} .
\]
Khususnya, karakter-karakter $\chi_1,\ldots, \chi_r$
membentuk basis ortogonal bagi subruang $R_k(G)$ dari $C(G)$.
\end{theorem}

\begin{proof}
Deskripsi $\chi_i$ diperoleh dari hasil-hasil pada bagian sebelumnya.
Keortogonalan mengikuti dari fakta bahwa $\psi_{i,\sigma}$
saling ortogonal.
\end{proof}
Perhatikan bahwa $\{\chi_1,\ldots,\chi_r\}$ hanya merupakan

basis \emph{ortogonal}; basis ini tidak selalu ortonormal.
Sebenarnya terdapat \emph{tiga} teras berbeda yang dapat dipertimbangkan

untuk setiap representasi tak tereduksi $V_i$.
Kita mempunyai teras biasa $\chi_i(g) \in k$ sebagai elemen dari
$\operatorname{End}_k(V_i)$.
Kita juga dapat mendefinisikan $\phi_i(g) \in F_i$ sebagai teras dari $g$ dalam
$\operatorname{End}_{F_i}(V_i)$.
Terakhir, kita mempunyai $\psi_i(g) \in F_i$ sebagai teras tereduksi
$\operatorname{Trd}(g)$, dengan memandang $g$ sebagai elemen aljabar sederhana sentral
$F_i$-$A_i$.
Semua teras ini terkait melalui
\[
\chi_i = \operatorname{Tr}_{F_i/k} \circ \phi_i, \quad
\phi_i = d_i \psi_i
\]
dengan mengingat bahwa $d_i$ adalah indeks Schur dari aljabar sederhana sentral
$F_i$-$A_i$ (secara ekuivalen, dari representasi $V_i$).
Terakhir, untuk setiap $\sigma \in X_i$, kita mempunyai $\psi_{i,\sigma} = \sigma
\circ \psi_i$.
Ini memberikan cara lain untuk memperoleh dekomposisi $V_i \otimes_k
\overline{k}$ di atas.
\subsection{Keterdefinisian Representasi}

Ingat dari bagian sebelumnya bahwa grup Galois absolut

mempermutasikan representasi-representasi tak tereduksi dari $\overline{k}G$.
Memang, jika $\rho_{i,\sigma} : G \to \operatorname{GL}(W_{i,\sigma})$
adalah pemetaannya, maka berlaku
\[
\tau \circ \rho_{i,\sigma} = \rho_{i,\tau(\sigma)}
\]
untuk $\tau \in \Gamma$.
Aksi ini sama dengan aksi yang diperoleh dengan bertindak pada
karakter-karakter
\[
\tau \circ \phi_{i,\sigma} = \phi_{i,\tau(\sigma)}
\]
dengan mengingat bahwa $\phi_{i,\sigma}: G \to \overline{k}$ hanyalah suatu
fungsi kelas tertentu.
\begin{theorem}

Misalkan $\overline{R}_k(G)$ adalah subhimpunan $R(G)$ yang terdiri atas karakter
yang fungsi kelasnya bernilai dalam $k$.
Maka $\overline{R}_k(G)$ mempunyai basis
\[
\frac{\chi_1}{d_1},
\frac{\chi_2}{d_2},
\cdots,
\frac{\chi_r}{d_r}.
\]
\end{theorem}
Jadi $\overline{R}_k(G)$ dapat diperoleh kembali sepenuhnya dari tabel karakter

atas $\mathbb{C}$. Tepatnya, $\chi_i/d_i$ dapat direkonstruksi
sebagai jumlah atas orbit-orbit $\Gamma$ dari karakter kompleks.
Untuk merekonstruksi $\overline{R}_k(G)$,
kita perlu mengetahui indeks-indeks Schur.

\begin{proposition}
Jika $G$ adalah grup abelian,
maka $R_k(G) = \overline{R}_k(G)$.
\end{proposition}

\begin{proof}
Jika $G$ abelian, maka $kG$ komutatif, sehingga semua aljabar pembagian
penyusunnya adalah medan. Jadi semua indeks Schur bernilai trivial.
\end{proof}

Kini kita siap membuktikan sebuah teorema penting dari Brauer, yang
pada dasarnya mereduksi studi keterdefinisian representasi atas
sembarang medan berkarakteristik $0$ ke kasus medan siklotomik.

\begin{theorem} \label{thm:define_over_cyclotomic}
Jika $G$ adalah grup hingga bereksponen $e$,
maka setiap representasi $G$ atas sembarang medan
berkarakteristik $0$
konjugat dengan suatu representasi yang terdefinisi atas
medan siklotomik $\mathbb{Q}(\zeta_e)$.
\end{theorem}

\begin{proof}
Misalkan $K=\mathbb{Q}(\zeta_e)$. Kita hendak menunjukkan bahwa $R_K(G)=R(G)$.
Menurut teorema Brauer, suatu elemen umum $\chi$ dari $R(G)$ dapat ditulis
dalam bentuk
\[
\rho = \sum_{i=1}^m c_i\operatorname{Ind}_{H_i}^G(\tau_i)
\]
dengan $c_i$ bilangan-bilangan bulat, $H_i$ subgrup-subgrup dari $G$,
dan setiap $\tau_i$ merupakan representasi berdimensi satu dalam $R(H_i)$.
Representasi berdimensi satu dari $H_i$ terdefinisi atas
$\mathbb{Q}(\zeta_e)$ karena eksponen $H_i$ paling besar $e$.
Dengan demikian, $\chi$ adalah kombinasi linear integral dari
representasi yang diinduksi dari $R_K(H_i)$, sehingga berada dalam
$R_K(G)$ sebagaimana diinginkan.

\end{proof}

\section{Ortogonalitas Kolom}
Kita telah melihat bahwa $\chi_1,\ldots,\chi_r$ membentuk basis ortogonal bagi
$R_k(G)$. Hal ini sama dengan menunjukkan bahwa \emph{baris-baris} tabel karakter
dari $G$ (atas $k$) saling ortogonal.
Namun, kini terdapat ``terlalu banyak'' kolom jika kita mengharapkan
tabel karakter berbentuk persegi.

Di sini kita menyelidiki persoalan tersebut.
Untuk setiap bilangan bulat positif $n$, definisikan \emph{operasi Adams}
$\Psi^n : C(G) \to C(G)$ sebagai fungsi yang memetakan
suatu fungsi kelas $\chi$ ke fungsi kelas $\Psi^n(\chi)$ yang didefinisikan oleh
\[
\Psi^n(\chi)(g) := \chi(g^n)
\]

untuk setiap $g \in G$.
\begin{proposition}
Setiap operasi Adams membatasi diri menjadi homomorfisme gelanggang
dari $R_k(G)$ ke dirinya sendiri.

\end{proposition}
\begin{proof}
Misalkan $\chi$ adalah karakter dari suatu representasi $(V,\rho)$.
Untuk suatu $g$, tinjau multihimpunan nilai eigen
\[
a_1,\ldots, a_m
\]
dari $\rho(g)$.
Jadi $\chi(g)=p_1$, dengan $p_1$ polinomial jumlah pangkat pertama
dalam $a_1,\ldots, a_m$.
Kita mempunyai $\chi(g^n)=p_n$, dengan $p_n$ polinomial jumlah pangkat
\[
a_1^n + \cdots + a_m^n.
\]
Menurut teorema dasar polinomial simetris, terdapat suatu
polinomial $f_n$ dengan koefisien bilangan bulat sedemikian sehingga
\[
p_n = f_n(e_1,\ldots,e_m)
\]
dengan $e_1,\ldots,e_m$ polinomial-polynomial simetris elementer

dalam $a_1,\ldots,a_n$.
Fungsi $f_n$ tidak bergantung pada nilai-nilai eigen, sehingga kita mempunyai
rumus
\[
\Psi^n(\chi) = f_n(\Lambda^1 \chi, \ldots, \Lambda^n \chi)
\]
dalam gelanggang $C(G)$. Karena pangkat eksterior suatu representasi adalah
representasi, $\Psi^n(\chi)$ berada dalam $R_k(G)$ sebagaimana diinginkan.

\end{proof}
Ingat bahwa \emph{eksponen} $e$ dari grup hingga adalah kelipatan persekutuan terkecil
dari orde semua elemen $g \in G$.
Secara ekuivalen, eksponen $e$ adalah bilangan bulat positif terkecil sedemikian sehingga

$g^e=1$ untuk semua $g \in G$.
\begin{lemma}
Jika $n$ relatif prima terhadap $e$,
maka fungsi $g \mapsto g^t$ merupakan bijeksi dari $G$ ke dirinya sendiri,
yang memetakan kelas konjugasi ke kelas konjugasi.

\end{lemma}
\begin{proof}
Misalkan $s$ adalah invers perkalian dari $n$ modulo $e$,
yang ada karena $t$ relatif prima terhadap $e$.
Fungsi $g \mapsto g^s$ adalah fungsi invers dari
$g \mapsto g^t$ karena $g^{ts}=g^1=g$.

\end{proof}

Dari sini segera diperoleh:
\begin{corollary}
Jika $t$ relatif prima terhadap $e$,
maka $\Psi^t : R(G) \to R(G)$ adalah suatu isomorfisme.

\end{corollary}
Jadi, untuk $t$ yang sesuai, operasi Adams $\Psi^t$

mempermutasikan \emph{kolom-kolom} tabel karakter dari $R(G)$.
Mengingat Teorema~\ref{thm:define_over_cyclotomic},
kita dapat mengasumsikan bahwa $k \subseteq K=\mathbb{Q}(\zeta_e)$
dengan $e$ eksponen dari $G$.
Misalkan $\Gamma$ adalah grup Galois dari $K/k$.
Perhatikan bahwa $\Gamma$ adalah subgrup dari
$(\mathbb{Z}/e\mathbb{Z})^\times$.
Misalkan $\Gamma$ dibangkitkan oleh $\sigma_1,\ldots,\sigma_s$
dengan $\sigma(\zeta_e)=\zeta_e^t$ untuk suatu $t_1,\ldots,t_s$ di

$(\mathbb{Z}/e\mathbb{Z})^\times$.
Kini kita mempunyai dua aksi $\Gamma$ pada $R(G)$.
Untuk $\gamma \in \Gamma$ dan $\chi \in R(G)$, kita mempunyai
\[
(\gamma \ast \chi)(g) := \gamma(\chi(g)) 
\]
untuk semua $g \in G$,
dengan memandang $\chi: G \to K$ sebagai fungsi berkodomain $K \subset
\overline{k}$.
Kita juga mempunyai aksi yang menggunakan operasi Adams,
yang dibangkitkan oleh
\[
\sigma_i \ast' \chi := \Psi^{t_i}(\chi)
\]

untuk $1 \le i \le s$.
\begin{proposition}
Kedua aksi $\Gamma$ pada $R(G)$ tersebut berimpit.

\end{proposition} 
\begin{proof}
Nilai eigen dari setiap matriks $M$ dalam setiap representasi $G$ merupakan
akar kesatuan ke-$e$.
Jadi, jika $\chi$ adalah karakter dalam $R(G)$, maka
\[
(\sigma_i \circ \chi)(g) = \Psi^{t_i}(\chi)(g)
\] 
untuk setiap $\sigma \in \Gamma$.

\end{proof}
Perhatikan bahwa kita dapat mendefinisikan aksi $\Gamma$ secara langsung pada kelas-kelas
konjugasi $G$ dengan menggunakan identifikasi di atas.
Berdasarkan hal ini, kita mendefinisikan \emph{kelas-$\Gamma$}
(atau \emph{kelas-$K$}) dari $G$ sebagai gabungan orbit-orbit $\Gamma$ dari
kelas-kelas konjugasi.

Dengan demikian kita memperoleh:
\begin{theorem}
$\overline{R}_k(G)$ tepat merupakan subhimpunan $R(G)$ yang karakter-karakter
terkaitnya tetap terhadap aksi $\Gamma$.

\end{theorem}

Sebagai akibat langsung, kita memperoleh:
\begin{corollary}
Jumlah karakter tak tereduksi dari $G$ atas $k$
sama dengan jumlah kelas-$\Gamma$ dari $G$.

\end{corollary}
Perhatikan bahwa aksi khusus $\Gamma$ pada $G$ ditentukan oleh suatu
pilihan isomorfisme antara $\Gamma$ dan citranya di dalam
$(\mathbb{Z}/e\mathbb{Z})^\times$, tetapi kelas-kelas $\Gamma$ tidak
bergantung pada pilihan ini.
Jadi kelas-kelas $\Gamma$ dari $G$ sepenuhnya bersifat teori-grup.
Khususnya, kita tidak perlu menghitung tabel karakter
untuk menentukan jumlah representasi tak tereduksi (seperti halnya atas

$\mathbb{C}$).
Namun, dari analisis ini kita hanya memperoleh $\overline{R}_k(G)$.

Masih diperlukan pekerjaan lebih lanjut untuk menentukan indeks-indeks Schur.

\subsection{Representasi atas bilangan rasional}
Di sini kita meninjau kasus khusus penting $k=\mathbb{Q}$.
Misalkan $G$ suatu grup hingga dengan eksponen $e$.
Seperti di atas, $K=\mathbb{Q}(\zeta_e)$ dengan $\zeta_e$ akar kesatuan primitif
ke-$e$.
Menurut hasil-hasil standar tentang medan siklotomik, grup Galois
$\Gamma = \operatorname{Gal}(K/k)$ isomorfik

dengan seluruh grup $(\mathbb{Z}/e\mathbb{Z})^\times$.

Dalam kasus ini, kelas-kelas $\Gamma$ di atas mempunyai deskripsi sederhana:
\begin{lemma}
Dua elemen $g,h \in G$ termasuk dalam kelas-$\mathbb{Q}$ yang sama jika dan
hanya jika keduanya membangkitkan subgrup siklik yang saling konjugat.

\end{lemma}
\begin{proof}
Perhatikan bahwa $\langle g \rangle = \langle h \rangle$ jika dan hanya jika
$g^s = h$ dan $h^t= g$ untuk suatu bilangan bulat $s$ dan $t$.
Perhatikan bahwa $t,s$ relatif prima terhadap $n=\langle g \rangle$.
Karena $n$ membagi eksponen $e$ dari $G$, kita mempunyai homomorfisme grup
surjektif $(\mathbb{Z}/e\mathbb{Z})^\times \to (\mathbb{Z}/n\mathbb{Z})^\times$.
Jadi kita dapat mengasumsikan $s$ dan $t$ relatif prima terhadap $e$.
Karena $\Gamma \cong (\mathbb{Z}/e\mathbb{Z})^\times$,
kita menyimpulkan bahwa $\langle g \rangle = \langle h \rangle$ jika dan hanya jika
$g$ dan $h$ berada dalam orbit $\Gamma$ yang sama.
Pernyataan ini tetap berlaku untuk kelas-kelas konjugasi. 
\end{proof}

Kita menyimpulkan teorema berikut:

\begin{theorem}
Jumlah kelas isomorfisme representasi tak tereduksi dari
grup hingga $G$ atas $\mathbb{Q}$ sama dengan jumlah kelas konjugasi
subgrup siklik dari $G$.
\end{theorem}

Kita menyoroti kesejajaran menarik berikut dengan representasi
permutasi dan representasi kompleks:

\begin{itemize}
\item Jumlah representasi \textbf{permutasi} tak tereduksi berkorespondensi
bijektif dengan jumlah kelas konjugasi \textbf{subgrup}.
\item Jumlah representasi \textbf{rasional} tak tereduksi berkorespondensi
bijektif dengan jumlah kelas konjugasi \textbf{subgrup
siklik}.
\item Jumlah representasi \textbf{kompleks} tak tereduksi berkorespondensi
bijektif dengan jumlah kelas konjugasi \textbf{elemen}.
\end{itemize}

\section{Representasi atas bilangan real}

Misalkan $V$ suatu representasi kompleks berdimensi hingga
dari suatu grup hingga $G$.

\begin{proposition} \label{prop:Ginner}
$V$ mempunyai hasil kali dalam yang invarian-$G$.
\end{proposition}

\begin{proof}
Misalkan $b : V \times V \to \mathbb{C}$ suatu hasil kali dalam (yang belum tentu
invarian-$G$).
Definisikan fungsi $B : V \times V \to \mathbb{C}$ dengan
\[
B(v,w) = \sum_{g \in G}b(gv,gw)
\]
untuk $v, w \in V$.
Perhatikan bahwa $B$ juga seskuilinear
dan bahwa $B(v,v) > 0$ untuk semua $v \ne 0$.
Jadi $B$ juga merupakan hasil kali dalam.
\end{proof}

Proposisi sebelumnya mempunyai interpretasi berikut:

\begin{corollary}
Setiap subgrup hingga dari $\operatorname{GL}_n(\mathbb{C})$
konjugat dengan suatu subgrup hingga dari grup uniter
$\operatorname{U}(n)$.
\end{corollary}

Yang penting, secara umum kita hanya dapat menghasilkan bentuk \emph{seskuilinear};
mungkin tidak ada bentuk \emph{bilinear} tak trivial.
Sebagai contoh, suatu representasi kompleks berdimensi satu yang tak trivial dari
$G=\mathbb{Z}/3\mathbb{Z}$ tidak mempunyai bentuk bilinear invarian-$G$ yang tak trivial;
jika ada, kita memperoleh kontradiksi
\[
1=(v,v)=(gv,gv)=(\zeta_3v,\zeta_3,v)=\zeta_3^2(v,v)=\zeta_3^2
\]
untuk vektor $v$ bernorma satu dan pembangkit $g$ dari $G$.

Namun, kita tahu bahwa bentuk seskuilinear dan bentuk bilinear
pada dasarnya sama atas medan tertutup real.
Fakta ini ternyata mengarah pada suatu teori yang menarik. 
Kita mengatakan bahwa $V$ \emph{dapat direalisasikan atas $\mathbb{R}$} jika terdapat
representasi real $W$ dari $G$ sedemikian sehingga $W \otimes_\mathbb{R}
\mathbb{C} \cong V$.

\begin{theorem}[Frobenius-Schur] \label{thm:FrobSchur}
Suatu representasi kompleks $V$ dapat direalisasikan atas $\mathbb{R}$ jika dan hanya
jika terdapat bentuk bilinear simetris tak degenerat yang invarian-$G$
pada $V$.
\end{theorem}

\begin{proof}
Pertama, andaikan $V$ dapat direalisasikan atas $\mathbb{R}$.
Misalkan $W$ representasi real sedemikian sehingga $V = W \otimes_\mathbb{R}
\mathbb{C}$. Perhatikan bahwa $W$ mempunyai bentuk bilinear simetris tak degenerat yang
invarian-$G$, yakni $q$, dengan argumen yang sama seperti Proposisi~\ref{prop:Ginner}
(hasil kali dalam sama dengan bentuk bilinear simetris tak degenerat
atas $\mathbb{R}$). Maka $q \otimes_\mathbb{R} \mathbb{C}$
adalah bentuk bilinear simetris tak degenerat yang invarian-$G$ pada $V$.

Sebaliknya, andaikan kita mempunyai bentuk bilinear simetris tak degenerat $B$
pada $V$. Kita juga mempunyai hasil kali dalam yang invarian-$G$, yakni $(-,-)$.
Kita mendefinisikan fungsi $\varphi : V \to V$ sedemikian sehingga
\[
B(x,y)=(\varphi(x),y)
\]
untuk semua $x,y \in V$
(fungsi ini tunggal dan terdefinisi dengan baik karena $B$ dan hasil kali dalam
keduanya tak degenerat).
Perhatikan bahwa $\varphi$ antilinear dan bijektif.
Jadi $\varphi^2$ merupakan automorfisme linear dari $V$.

Kita mengklaim bahwa $\varphi^2$ adalah operator Hermite definit positif.
Memang,
\[
(\varphi^2(x),y)=B(\varphi(x),y)=B(y,\varphi(x))=(\varphi(y),\varphi(x))
\]
untuk semua $x,y \in V$.
Dengan demikian,
\[
\overline{(\varphi(x),\varphi(y))}=\overline{(\varphi^2(y),x)}.
\]
Karena hasil kali dalam bersifat simetris-konjugat, kita memperoleh
\[
(\varphi^2(x),y)=(x,\varphi^2(y))
\]
dan menyimpulkan bahwa $\varphi^2$ bersifat Hermite.
Selain itu, $(\varphi^2(x),x)=(\varphi(x),\varphi(x))$
menunjukkan bahwa operator tersebut definit positif.

Setiap operator Hermite definit positif $\varphi^2$ mempunyai akar kuadrat
Hermite definit positif tunggal $\omega$ sedemikian sehingga
$\omega^2=\varphi^2$. Selain itu, $\omega$ komutatif dengan $\varphi$.
Misalkan $\sigma = \varphi \omega^{-1}$.
Kita melihat bahwa $\sigma : V \to V$ adalah pemetaan antilinear dengan kuadrat
sama dengan identitas.
Selain itu, $\sigma$ ekuivarian-$G$.

Karena $\sigma$ antilinear, khususnya ia
linear-$\mathbb{R}$. 
Misalkan $V_R$ adalah ruang eigen-$1$ (real) dari $\sigma$, dan $V_I$ adalah
ruang eigen-$-1$ (real) dari $\sigma$. Karena $\sigma$ antilinear,
kita memperoleh $V_I=iV_R$. Dengan demikian, $V=V_R \oplus iV_I$
sebagaimana diinginkan.
\end{proof}

Sekarang andaikan $V$ tak tereduksi dan misalkan $\chi$ karakternya.

\begin{definition}
\emph{Indikator Frobenius--Schur} dari $V$ adalah bilangan
\[
\iota(V) := \frac{1}{|G|}\sum_{g \in G} \chi(g^2).
\]
\end{definition}

Untuk memahami teori representasi atas $\mathbb{R}$,
perlu dicatat bahwa hanya ada tiga aljabar pembagian
atas $\mathbb{R}$:
medan real $\mathbb{R}$,
medan kompleks $\mathbb{C}$, dan
kuaternion Hamilton $\mathbb{H}$.
Representasi tak tereduksi $V$ merupakan modul sederhana dari
$A \otimes_{\mathbb{R}} \mathbb{C}$, dengan $A$ subaljabar sederhana dari
$\mathbb{R}G$ dan $A \cong \operatorname{M}_n(D)$, dengan $D$ salah satu
dari ketiga kasus di atas.

Jadi kita memperoleh trikotomi perilaku real dari representasi
kompleks berdasarkan aljabar pembagian yang
bersesuaian dengannya.

\pagebreak

\noindent
\textbf{Kompleks:}
\begin{itemize}
\item Aljabar pembagian yang bersesuaian dengan $V$ adalah $\mathbb{C}$,
\item indikator Frobenius--Schur adalah $\iota(V)=0$,
\item $V$ tidak mempunyai bentuk bilinear tak degenerat yang invarian-$G$,
\item $\chi$ tidak bernilai real,
\item $V$ tidak dapat direalisasikan atas $\mathbb{R}$, dan
\item $V$ tidak isomorfik dengan dualnya $V^\vee$.
\end{itemize}

\noindent
\textbf{Real:}
\begin{itemize}
\item Aljabar pembagian yang bersesuaian dengan $V$ adalah $\mathbb{R}$,
\item indikator Frobenius--Schur adalah $\iota(V)=1$,
\item $V$ mempunyai bentuk bilinear simetris tak degenerat yang invarian-$G$,
\item $\chi$ bernilai real,
\item $V$ dapat direalisasikan atas $\mathbb{R}$, dan
\item $V \cong V^\vee$.
\end{itemize}

\noindent
\textbf{Kuaternionik:}
\begin{itemize}
\item Aljabar pembagian yang bersesuaian dengan $V$ adalah $\mathbb{H}$,
\item indikator Frobenius--Schur adalah $\iota(V)=-1$,
\item $V$ mempunyai bentuk bilinear antisimetris tak degenerat yang invarian-$G$,
\item $\chi$ bernilai real,
\item $V$ tidak dapat direalisasikan atas $\mathbb{R}$, tetapi
\item $V \oplus V$ dapat direalisasikan atas $\mathbb{R}$, dan
\item $V \cong V^\vee$.
\end{itemize}

Kita hanya menguraikan garis besar alasan karakterisasi di atas berlaku,
dan menyerahkan rinciannya kepada pembaca (atau \cite[\S{13.2}]{Serre}).
Sebagian besar karakterisasi segera mengikuti teori Wedderburn,
sehingga kita berfokus pada bentuk-bentuk bilinear.

Perhatikan bahwa bentuk bilinear ekuivalen dengan pemetaan antara $V$ dan
$V^\vee$, sedangkan ketaktrivialan ekuivalen dengan ketakdegeneratan menurut
Lema Schur karena kita menghendaki keinvarianan-$G$.
Lema Schur juga mengatakan bahwa bentuk bilinear yang invarian-$G$
tunggal hingga perkalian dengan skalar.
Dekomposisi bentuk bilinear menjadi bagian simetris dan bagian antisimetris
bersifat invarian-$G$, sehingga ketunggalan hingga skalar menyiratkan bahwa bentuk itu
harus simetris atau antisimetris.
Keberadaan bentuk bilinear \emph{simetris} ekuivalen dengan keterrealisasian atas
bilangan real menurut Teorema~\ref{thm:FrobSchur} di atas.

Terakhir, kita meninjau indikator Frobenius--Schur.
Suatu bentuk bilinear simetris yang invarian-$G$ pada $V$ ekuivalen dengan
keberadaan subrepresentasi trivial dari $\operatorname{Sym}^2(V)$.
Demikian pula, bentuk bilinear antisimetris yang invarian-$G$ pada $V$
ekuivalen dengan keberadaan subrepresentasi trivial dari
$\operatorname{Alt}^2(V)$.
Ingat kembali rumus untuk karakter dari
kuadrat simetris dan kuadrat berselang-seling:
\[
\chi_{\operatorname{Sym}^2(V)} = \chi_V^2 + \Psi^2(\chi_V)
\textrm{ dan }
\chi_{\operatorname{Alt}^2(V)} = \chi_V^2 - \Psi^2(\chi_V).
\]
Jadi, $(1,\chi_{\operatorname{Sym}^2(V)})=1$ jika dan hanya jika $V$ real,
$(1,\chi_{\operatorname{Alt}^2(V)})=1$ jika dan hanya jika $V$ kuaternionik,
dan keduanya bernilai $0$ dalam kasus lainnya.

Sekarang cukup kita perhatikan bahwa
\[
\chi_V^2 = \chi_{\operatorname{Sym}^2(V)} + \chi_{\operatorname{Alt}^2(V)}
\]
dan
\[
\iota(V) := \frac{1}{|G|}\sum_{g \in G} \chi(g^2)
=(1, \Psi^2(\chi_V)).
\]

\bibliographystyle{alpha}
\bibliography{rep_theory}

\end{document}


